% Source-derived chapter 07: Exponential and logarithm
% Original source: chabauty-coleman-bound-surfaces
\chapter{Exponential and logarithm}
\label{chabauty-coleman-bound-surfaces-chapter-07}
\source{chabauty-coleman-bound-surfaces}

\begin{remark}[Source section]
\label{chabauty-coleman-bound-surfaces-chapter-07-source}
\source{chabauty-coleman-bound-surfaces}
\source{chabauty-coleman-bound-surfaces}
This chapter reproduces the corresponding section of the retrieved
LaTeX source, including its definitions, statements, examples, and
proofs. It has not yet been translated into Lean.
\notready
\end{remark}

% The source section title was: Exponential and logarithm
\label{SecLogExp}
\source{chabauty-coleman-bound-surfaces}

Let us keep the notation from Section \ref{NotationLoc}. Let $\Acal$ be an abelian variety over $\Spec R$ of relative dimension $n\ge 1$. 
%%%%
%%%%
%%%%
\subsection{Local parameters} \label{SecLogExp1}
\source{chabauty-coleman-bound-surfaces}


Let $\sigma:\Spec R\to \Acal$ be a section with image $\Zcal=\sigma(\Spec R)$. The closed point of $\Zcal$ is $x_0=\sigma(\pfrak)$. Let $\widehat{\Ocal}_{\Acal ,x_0,\Zcal}$ be the completion of $\Ocal_{\Acal,x_0}$ along $\Zcal$.  That is, if $\Ical_\Zcal$ is the ideal sheaf of $\Zcal$, then $\widehat{\Ocal}_{\Acal ,x_0,\Zcal}$ is the completion of $\Ocal_{\Acal,x_0}$ with respect to the ideal $\Ical_{\Zcal,x_0}\subseteq\Ocal_{\Acal,x_0}$.  

 Let $y_1,...,y_n$ be variables. We say that elements $t_1,...,t_n\in \mfrak_{\Acal,x_0}$  are \emph{$R$-local parameters} along $\sigma$ (or $\Zcal$) if the rule $y_j\mapsto t_j$ induces a continuous isomorphism of $R$-algebras $\widehat{\Ocal}_{\Acal,x_0,\Zcal}\simeq R[[y_1,...,y_n]]$.
%%
%%
\begin{lemma}[Choice of local parameters]
\source{chabauty-coleman-bound-surfaces}
\notready\label{LemmaChoiceLocParam} With the previous notation, let $\Xcal$ be a regular closed subscheme of $\Acal$ containing $\Zcal$ which is smooth of relative dimension $d\ge 0$ over $\Spec R$.  There are $t_j\in \mfrak_{\Acal,x_0}$ for $1\le j\le n$ such that $t_1,...,t_n$ are $R$-local parameters along $\Zcal$ and $t_1,...,t_{n-d}$ are a system of local equations for $\Xcal$ in $\Acal$ along $\Zcal$. In particular, $R$-local parameters along $\Zcal$ exist.
\source{chabauty-coleman-bound-surfaces}
\end{lemma}
\begin{proof} By \cite{SGA1} Expos\'e II, Th\'eor\`eme 4.15 with $X=\Acal$, $Y=\Xcal$, $S=\Spec R$, and $x=x_0$  we get the existence of  a regular sequence $t_1,...,t_{n-d}\in \mfrak_{\Acal,x_0}$ that generates the ideal $\Ical_{\Xcal, x_0}$ of $\Xcal$ in $\Ocal_{\Acal, x_0}$. Applying \cite{SGA1} Expos\'e II, Corollaire 4.17 with $X=\Acal$, $Y=\Spec R$, $i=\sigma$, and $y=\pfrak\in \Spec R$ we get that  $\widehat{\Ocal}_{\Acal,x_0,\Zcal}$ is a power series ring over $R$ in $n$ variables. Analyzing the proof, one sees that the variables of this power series ring can be constructed by extending any regular sequence in $\Ical_{\Zcal, x_0}$, see \cite{SGA1} Expos\'e II, Remarques  4.14. We conclude since $\Ical_{\Xcal,x_0}\subseteq \Ical_{\Zcal,x_0}$.
\end{proof}
%%
%%
Let $\ttt=(t_1,...,t_n)$ be a choice of $R$-local parameters along the identity section $e$. These $R$-local parameters determine a formal group $\Fcal=(\Fcal_1,...,\Fcal_n)$ with $\Fcal_j\in R[[X_1,...,X_n,Y_1,...,Y_n]]$ for each $1\le j\le n$, where the $X_j$ and $Y_j$ are variables, see \cite{SilvermanHindry} Lemma C.2.4. 

Given an integer $m$, let $\Psi^{[m]}=(\Psi^{[m]}_1,...,\Psi^{[m]}_n)$ be the power series expansion of the morphism $[m]:\Acal\to \Acal$ of multiplication by $m$ in terms of $\ttt$. Then, integrality of the formal group $\Fcal$ gives
%%
%%
\begin{lemma}
\source{chabauty-coleman-bound-surfaces}
\notready\label{LemmaPsiR} $\Psi^{[m]}_j\in R[[\ttt]]$ for each $j=1,...,n$. 
\source{chabauty-coleman-bound-surfaces}
\end{lemma}
%%
%%

Although the notation does not explicitly indicate so, the coefficients of the power series $\Psi^{[m]}_j$  depend upon the choice of local parameters $\ttt$. Nevertheless, since $[1]=\Id$ we observe:
%%
%%
\begin{lemma}
\source{chabauty-coleman-bound-surfaces}
\notready\label{LemmaLinearPart} $\Psi^{[1]}_j =t_j$ for each $j=1,...,n$.
\source{chabauty-coleman-bound-surfaces}
\end{lemma}
%%
%%
  From \cite{Bourbaki} III.5.3 Proposition 2(i) we get
%%
%%
\begin{lemma}
\source{chabauty-coleman-bound-surfaces}
\notready\label{LemmaVanishPsi} The terms of degree less than $m$ in $\Psi^{[m]}_j$ vanish for each $j=1,...,n$.
\source{chabauty-coleman-bound-surfaces}
\end{lemma}
%%
%%

%%%%
%%%%
%%%%
\subsection{Power series construction} \label{SecLogExp2} Let $\psi^{[m]}_j$ be the homogeneous part of degree $m$ in $\Psi_j^{[m]}$. Thus, $\psi_j^{[1]} = t_j$ by Lemma \ref{LemmaLinearPart}. Let $x_1,...,x_n$ and $y_1,...,y_n$ be variables and for each $j=1,...,n$ we define 
\source{chabauty-coleman-bound-surfaces}
$$
\begin{aligned}
\Log^{\ttt}_j(\yy) & = \sum_{m=1}^\infty \frac{(-1)^m}{m}\cdot \Psi^{[m]}_j(\yy)\in K[[\yy]]\\
\Exp^{\ttt}_j(\xx) & = \sum_{m=1}^\infty \frac{1}{m!}\cdot \psi^{[m]}_j(\xx)\in K[[\xx]].
\end{aligned}
$$
These power series depend upon the choice of local parameters $\ttt$ as the notation indicates. When this choice is clear, we will simply write
$$
\Log^\ttt_j(\yy)=\sum_\alpha b_{j,\alpha}\cdot \yy^\alpha \quad \mbox{ and }\quad \Exp^\ttt_j(\xx)=\sum_\alpha c_{j,\alpha}\cdot \xx^\alpha.
$$
The coefficients $b_{j,\alpha}$ and $c_{j,\alpha}$ also depend on the choice of $\ttt$, although the notation is not explicit.

By \cite{Bourbaki} III.5.4 Prop. 3, these are the power series expansions of the exponential and logarithm maps of the $p$-adic Lie group $A(K)$ at the identity point $e$, with respect to the choice of local parameters $\ttt$. They converge on some $p$-adic neighborhood of $e$, but for our purposes we need a more precise discussion on convergence.


%%
%%
\begin{lemma}[Convergence of  logarithm]\label{LemmaCvLog} For each $j=1,...,n$, the radius of convergence of $\Log^{\ttt}_j$ is at least $1$. Furthermore, for each $\alpha\in \N^n$ we have $|b_{j,\alpha}|\le \|\alpha\|^{[K:\Q_p]
\source{chabauty-coleman-bound-surfaces}
\notready}$.
\source{chabauty-coleman-bound-surfaces}
\end{lemma}
\begin{proof} From Lemmas \ref{LemmaPsiR} and \ref{LemmaVanishPsi} we get $|b_{j,\alpha}|\le \max\{ |m^{-1}|: 1\le m\le\|\alpha\|\} \le \|\alpha\|^{[K:\Q_p]}$. The claim on the radius of convergence follows.
\end{proof}
%%
%%


%%
%%
\begin{lemma}[Convergence of the exponential] \label{LemmaCvExp} For each $j=1,...,n$ and each $\alpha\in \N^n$, we have $m!\cdot c_{j,\alpha}\in R$ where $m=\|\alpha\|$. Furthermore, $|c_{j,\alpha}|\le p^{[K:\Q_p](m-1)/(p-1)}$. In particular, the radius of convergence of $\Exp^{\ttt}_j$ is at least $p^{-[K:\Q_p]
\source{chabauty-coleman-bound-surfaces}
\notready/(p-1)}$.
\source{chabauty-coleman-bound-surfaces}
\end{lemma}
\begin{proof} Since $c_{j,\alpha}$ is a coefficient of $m!^{-1} \psi_j^{[m]}(\xx)$, we get $m!\cdot c_{j,\alpha}\in R$ from Lemma \ref{LemmaPsiR}. Letting $s_p(m)$ be the sum of the digits of $m$ in base $p$, we get
$$
|c_{j,\alpha}|\le \frac{1}{|m!|} =p^{[K:\Q_p](m-s_p(m))/(p-1)}\le p^{[K:\Q_p](m-1)/(p-1)}.
$$
The claim on the radius of convergence follows.
\end{proof}
%%
%%

%%%%
%%%%
%%%%
\subsection{Local linearization}\label{SecLocLin} We consider the open set $V = \pfrak\times ...\times \pfrak\subseteq K^n$ and the additive Lie group $\Tcal$ given by $V$ with the usual addition. The variables $\xx=(x_1,...,x_n)$ and $\yy=(y_1,...,y_n)$ will be considered as coordinates on $\Tcal$ and $V$ respectively.
\source{chabauty-coleman-bound-surfaces}

Since $\Acal\to \Spec R$ is proper, we have the reduction map $\red : A(K)=\Acal(K)\to A'(\kk)$ where $A'=\Acal\otimes \kk$. This is a group morphism. Then $U_e=\ker(\red)$ is a Lie subgroup of $A(K)$. 

The $R$-local parameters $\ttt=(t_1,...,t_n)$ define functions $t_j:U_e\to \pfrak$ which determine a $K$-analytic local chart $\chi^\ttt:U_e\to V$. The map $\chi^\ttt$ is bijective, although it does not need to be a group morphism.

%%
%%
\begin{lemma}
\source{chabauty-coleman-bound-surfaces}
\notready\label{LemmaKeyLogExp} Suppose that $p>\max\{\efrak +1 , \exp(\efrak/\exp(1))\}$. Then we have the following:
\source{chabauty-coleman-bound-surfaces}
\begin{itemize}
\item[(i)]  $\Log^\ttt_j$ and $\Exp^\ttt_j$ have radius of convergence larger than $1/q$ for each $j=1,...,n$. 
\item[(ii)] The power series $\Log^\ttt(\yy)=(\Log^\ttt_j(\xx))_{j=1}^n$ and $\Exp^\ttt(\xx)=(\Exp^\ttt_j(\yy))_{j=1}^n$ give analytic maps
$$
\Log^\ttt : V\to \Tcal\quad  \mbox{ and }\quad \Exp^\ttt:\Tcal\to V
$$
which are inverse to each other.
\item[(iii)] The maps $\widetilde{\Log}^\ttt$ and $\widetilde{\Exp}^\ttt$ defined by
$$
\widetilde{\Log}^\ttt=\Log^\ttt \circ \chi^\ttt: U_e\to \Tcal\quad  \mbox{ and }\quad \widetilde{\Exp}^\ttt=(\chi^\ttt)^{-1}\circ \Exp^\ttt:\Tcal\to U_e
$$
are isomorphisms of Lie groups, inverse to each other.
\end{itemize}
\end{lemma}
\begin{proof} (i)  Lemma \ref{LemmaCvLog} implies the result for $\Log^\ttt_j$. The condition $p>\efrak+1$ gives $[K:\Q_p]<(p-1)\ffrak$, hence $p^{[K:\Q_p]/(p-1)}<q$. Lemma \ref{LemmaCvExp} shows that $\Exp^\ttt_j$ has radius of convergence at least $p^{-[K:\Q_p]/(p-1)}$, which is larger than $1/q$.

(ii) As open sets of $K^n$, we have $V=\Tcal=B_n[1/q]$. Thus, (i) gives that $\Log^\ttt$ and $\Exp^\ttt$ are analytic on the domains $V$ and $\Tcal$ respectively, as maps to $K^n$.

Let us show that $\Log^\ttt$ maps $V$ to $\Tcal$. Since $B_n[1/q]=B_n(1)$ in $K^n$, it suffices to check that for all $\vv\in V$ and $j=1,...,n$ we have $|\Log^\ttt_j(\vv)|<1$. By Lemma \ref{LemmaCvLog} we have 
$$
|\Log^\ttt_j(\vv)|\le \max_{m\ge 1} m^{[K:\Q_p]}/q^m =\max\left\{1/q , \max_{m\ge 2} m^{[K:\Q_p]}/q^m\right\}. 
$$
By $p>\exp(\efrak/\exp(1))$ we have $m/\log m> \exp(1)> \efrak/\log p$ for all $m\ge 2$. Thus, $[K:\Q_p]\log m =\efrak\ffrak\log m < \ffrak m\log p=m\log q$. Hence, $m^{[K:\Q_p]}/q^m<1$ for each $m\ge 2$, proving  $|\Log^\ttt(\vv)_j|<1$.

To show that $\Exp^\ttt$ maps $\Tcal$ to $V$, it suffices to check that for each $\vv\in \Tcal$ and $j=1,...,n$ we have $|\Exp^\ttt_j(\vv)|<1$. Lemma \ref{LemmaCvExp} gives 
$$
|\Exp^\ttt_j(\vv)|\le \sup_{m\ge 1} p^{[K:\Q_p]\frac{m-1}{p-1}}/q^m = \sup_{m\ge 1} p^{[K:\Q_p]\frac{m-1}{p-1} -\ffrak m}.
$$
Since $p>\efrak+1$, we get $(m-1)\efrak < (p-1)m$. This gives $[K:\Q_p]\frac{m-1}{p-1} -\ffrak m<0$, proving $|\Exp^\ttt_j(\vv)|<1$.

The fact that $\Log^\ttt$ and $\Exp^\ttt$ are inverse to each other is a power series identity, see \cite{Bourbaki} III 5.4.

Finally, (iii) also follows from \cite{Bourbaki} III 5.4, where a neighborhood of $e$ is identified with an open set of $K^n$ by the choice of local parameters made in \cite{Bourbaki} III 5.3.
\end{proof}
%%
%%



%%%%%%%%%%%%%%%%%%%%%%%%%%%%%%%%%%%%%%
%%%%%%%%%%%%%%%%%%%%%%%%%%%%%%%%%%%%%%
%%%%%%%%%%%%%%%%%%%%%%%%%%%%%%%%%%%%%%
%%%%%%%%%%%%%%%%%%%%%%%%%%%%%%%%%%%%%%
%%%%%%%%%%%%%%%%%%%%%%%%%%%%%%%%%%%%%%
%%%%%%%%%%%%%%%%%%%%%%%%%%%%%%%%%%%%%%
